\section{Matrices}
Following our example from the end of the previous section, let's look again at the result of applying $T$ from \autoref{eq:simple_LT_example_single_2d_vec} on the vector $\vec{a}=\clrcolvecxy{3}{2}$.

Since we know how $T$ transforms the basis vectors $\eb{1}$ and $\eb{2}$ (\autoref{eq:simple LT example transformed basis vectors}), we can $T(\vec{a})$ in the following way:
\begin{equation}
	T(\vec{a}) = \colvec{3;4} = \overset{\substack{\LCa{a_{1}} \\ \downarrow}}{\LCa{3}}\underset{\substack{\uparrow \\ T(\eb{1})}}{\colvec{0;1}} + \overset{\substack{\LCb{a_{2}} \\ \downarrow}}{\LCb{2}}\underset{\substack{\uparrow \\ T(\eb{2})}}{\colvec{-2;0}} = \colvec{0\cdot\LCa{3}-2\cdot\LCb{2};1\cdot\LCa{3}+0\cdot\LCb{2}}.
	\label{eq:developing matrix notation step 1}
\end{equation}

Then, we collect all the coefficients in black into a single structure, which we call a \newdef{matrix}:
\begin{equation}
	M =
	\begin{bmatrix}
		0 & -2 \\
		1 & 0
	\end{bmatrix}.
	\label{eq:developing matrix notation step 2}
\end{equation}

Finally, we define the action of the matrix $M$ on the vector $\vec{a}$ as following:
\begin{equation}
	M\vec{a} = 
	\begin{bmatrix}
		0 & -2 \\
		1 & 0
	\end{bmatrix}\clrcolvecxy{3}{2} = \colvec{0\cdot\LCa{3}-2\cdot\LCb{2};1\cdot\LCa{3}+0\cdot\LCb{2}} = T(\vec{a}).
	\label{eq:developing matrix notation step 3}
\end{equation}

Of course, we should make sure that the matrix $M$ we just defined works for other vectors as well. If we apply $M$ to the second vector from our example, $\vec{b}=\clrcolvecxy{2}{-4}$, we get
\begin{equation}
	M\vec{b} = 
	\begin{bmatrix}
		0 & -2 \\
		1 & 0
	\end{bmatrix}\clrcolvecxy{2}{-4} = \colvec{0\cdot\LCa{2}-2\cdot\LCb{-4};1\cdot\LCa{2}+0\cdot\LCb{-4}} = \colvec{8;2} = T(\vec{b}),
	\label{eq:validating matrix works for b}
\end{equation}
which is exactly the result we got in \autoref{eq:second_simple_LT_example_single_2d_vec}. Great!

Why not continue checking how this matrix-vector product we developed works for other vectors? First, let's check $\eb{1}$ and $\eb{2}$ for which we already know the results (\autoref{eq:simple_LT_example_single_2d_vec}):
\begin{equation}
	\begin{aligned}
		M\eb{1} &= 
		\begin{bmatrix}
			0 & -2 \\
			1 & 0
		\end{bmatrix}\clrcolvecxy{1}{0} = \colvec{0\cdot\LCa{1}-2\cdot\LCb{0};1\cdot\LCa{1}+0\cdot\LCb{0}} = \colvec{0;1} = T(\eb{1}),\\
		M\eb{2} &= 
		\begin{bmatrix}
			0 & -2 \\
			1 & 0
		\end{bmatrix}\clrcolvecxy{0}{1} = \colvec{0\cdot\LCa{0}-2\cdot\LCb{1};1\cdot\LCa{0}+0\cdot\LCb{1}} = \colvec{-2;0} = T(\eb{2}),
		\label{eq:}
	\end{aligned}
\end{equation}
exactly the results we expected.

Perhaps you've noticed something by now - the columns of $M$ are exactly the vectors $T(\eb{1})$ and $T(\eb{2})$ in the right order:
% \begin{equation}
% 	M = 
% 	\begin{bNiceMatrix}[margin]
% 		\tikzmark{e1}\Block[fill=xpurple!25, rounded-corners]{2-1}{}0 & \tikzmark{e2}\Block[fill=xorange!25, rounded-corners]{2-1}{}-2 \\
% 		1 & 0
% 		\begin{tikz}[overlay, remember picture] {\node[above of=e1] {$\eb{1}$}; \node[above of=e2] {$\eb{2}$}}
% 	\end{bNiceMatrix}.
% 	\label{eq:}
% \end{equation}

\begin{equation*}
	M = \begin{tikzpicture}[baseline=-0.5ex]
		\matrix [matrix of math nodes,
		nodes={rectangle, minimum size=15pt}, % should somehow also make a vertical space to the bottom of the matrix
		inner sep=0pt,
		column sep=8pt,
		left delimiter={[}, right delimiter={]},
		] (M)
		{
			0 & -2 \\
			1 & 0 \\
		};
		\begin{scope}[on background layer]
			\filldraw[xpurple, rounded corners, opacity=0.25] (M-1-1.north west) rectangle (M-2-1.south east);
			\filldraw[xorange, rounded corners, opacity=0.25] (M-1-2.north west) rectangle (M-2-2.south east);
		\end{scope}

		\tiny
		\node[xpurple, above of=M-1-1] (eb1text) {$T(\eb{1})$};
		\draw[-stealth, xpurple] (eb1text) -- (M-1-1.north);
		\node[xorange, above of=M-1-2] (eb2text) {$T(\eb{2})$};
		\draw[-stealth, xorange] (eb2text) -- (M-1-2.north);
\end{tikzpicture}.
	\label{eq:}
\end{equation*}
